% B-58, 15.09.2026: F2 mit Realisierungs- und Evaluationsbeleg verbunden.
% P10-4a Grafik v03, 14.09.2026: Durchlauf, GitHub-Ernte, Quellenweg und Beziehungserzeugung integriert.
% P10-4a, 14.09.2026: Mechanismus, Quellenreichweite und v4-Nachweis angeglichen.
% M17, 12.09.2026: Initialisierung in vier erklärende Absätze gegliedert; Schreibfolge am Code gegengeprüft. Betriebspfad unverändert.
% M07, 11.09.2026: F1-Anschlüsse gebündelt; Originalsprecherin und
% gespeichertes Präfix getrennt. Keine Änderung historischer Artefakte.
% Quellen: meeting-4-ux-block-l.txt:33, F1 consumable.json (visit-planning-
% announcement), human-decisions.json und github-forward-apply-report.json.
% Bootstrap- und spätere Fortschreibungsbelege bleiben getrennt.

\section{Zwei repräsentative End-to-End-Pfade}
\label{sec:endtoend-pfade}

Die beiden Pfade führen die zuvor erläuterten Architekturstationen
zusammen: Zunächst wird die Herstellung des ersten Projektbestands
beschrieben, anschließend die Fortschreibung am Besuchsbeispiel aus
Abschnitt~\ref{sec:zielbild-anforderungen-systemgrenze}. Die Belege
stammen aus unterschiedlichen Ausführungskontexten; sie bilden
keinen gemeinsamen End-to-End-Lauf.

\subsection*{Initiale Herstellung: vom leeren Projekt zum ersten Bestand}

Bei der Initialisierung ist noch kein Core angelegt, gegen den
neue Informationen eingeordnet werden könnten. Der Pfad lässt sich
in vier Abschnitte gliedern:

\emph{Erstens: Evidenz gewinnen.} Aus dem Transkript entsteht die
verwendbare Evidenzsicht
(Abschnitt~\ref{sec:evidenzgebundene-verarbeitung}). Die menschliche
Adjudikation ist selektiv: Vorgelegt werden beanstandete Claims und
Hinweise auf ungenutzte Quellstellen. Bei leerer Vorlagenliste
entfällt der Entscheidungspunkt.

\emph{Zweitens: fachliche Artefakte ableiten.} Modellgestützte
Ableitungen erzeugen aus dieser Evidenzsicht Anforderungs-, Fragen-
und Architektursichten
(Abschnitt~\ref{sec:saeule-semantische-transformation}). E44-06
illustriert die Herkunftskette vom Claim zum abgeleiteten Requirement.

\emph{Drittens: den Core initialisieren.} Der deterministische
Erst-Seed überführt die abgeleiteten Inhalte in den ersten
Projektbestand. Er schreibt die Anforderungs- und Architekturelemente
über die gemeinsame Speichernaht \emph{ohne gesonderte menschliche
Einzelautorisierung}. Existiert bereits ein Core, bricht dieser
Schritt ab. Der Überschreibschutz und die vorgelagerte selektive
Evidenz-Adjudikation begrenzen den Vorgang, ersetzen aber keine
Einzelautorisierung dieser Elemente.

\emph{Viertens: den Bestand strukturieren.} Die anschließende
Themenstrukturierung und Arbeitspaketbildung ergänzen den bereits
angelegten Bestand. Ihre Vorschläge durchlaufen die vorgesehenen
menschlichen Gates (Abschnitt~\ref{sec:governance-mutation}). Der
erste Schreibvorgang ist zu diesem Zeitpunkt somit bereits erfolgt.

Das in Kapitel~\ref{chap:explorative-problemklaerung} hergeleitete
Autorisierungsziel ist im Initialisierungspfad nur teilweise
realisiert. Kapitel~\ref{chap:evaluation} begrenzt entsprechend die
Reichweite der Governance-Nachweise. Dokumentierte Bootstrap-Läufe
sind im Evaluationsanhang zugeordnet.

\subsection*{Laufender Betrieb: ein Inhalt durch alle Stationen}

Im synthetischen Folgegespräch des tatsächlich ausgeführten Falls F1
wird die Besuchsankündigung als neues Thema in einen vorhandenen
Projektbestand eingebracht. Tabelle~\ref{t:beispiel-besuchsplanung}
verfolgt diesen Inhalt bis zu den real erzeugten GitHub-Issues.
Die Einordnung des Falls und seiner Dokumentationsgrenzen erfolgt
in Abschnitt~\ref{sec:eval-fortschreibung}.

% P10-4a, 14.09.2026: Mechanismus, Quellenreichweite und v4-Nachweis angeglichen.
% M07 / B-20, 11.09.2026: Quellenzeile aus meeting-4-ux-block-l.txt:33;
% F1-Ledger-Beleg hat abweichendes Präfix Leitung. Nicht stillschweigend
% korrigierter Laufbestand. Tabelleninhalte im historischen F1-Stand,
% keine späteren Erweiterungen aus dem heutigen Core übernehmen.

\begin{table}[htbp]
  \centering
  \small
  \begin{tabular}{>{\raggedright\arraybackslash}p{0.20\textwidth}>{\raggedright\arraybackslash}p{0.72\textwidth}}
    \hline
    \textbf{Station} & \textbf{Inhalt im dokumentierten Lauf} \\
    \hline
    Quelle & Aussage der Angehörigen-Vertreterin im Originaltranskript: „Angehörige sollen Besuche bei ihrer Bezugsperson vorab in der App ankündigen können, mit Datum und Uhrzeit, und die Pflegenden der Einrichtung sollen eine Übersicht der angekündigten Besuche sehen.`` \\
    Ledger & Ein quellgebundener Claim bündelt beide Teilaussagen („Angehörige sollen \ldots{} ankündigen können, und Pflegende \ldots{} sollen eine Übersicht der angekündigten Besuche sehen``) und enthält einen gespeicherten Belegtext; dessen Sprecherpräfix weicht vom Original ab (siehe Fließtext). \\
    Fachliche Ableitung & Die modellgestützte Ableitung zerlegt den gebündelten Claim in zwei normative Anforderungs-Aussagen mit vereinheitlichter Modalität („Angehörige \emph{müssen} \ldots{} ankündigen können`` / „Pflegende \emph{müssen} eine Übersicht \ldots{} sehen können``); der Claim-Bezug bleibt als Herkunftsangabe erhalten. \\
    Entscheidungsvorlage & Der Einordnungsplan stellt zwölf Operationen zur Entscheidung; für die Besuchsankündigung eine Neuaufnahme mit Begründung („Für Besuchsankündigungen existiert im Core keine passende Requirement-Entität``) samt Claim-Referenz. \\
    Übernahme & Der Mensch übernahm am Ingest-Gate acht Operationen und wies vier mit protokollierter Begründung ab; die deterministische Anwendung vergab die Core-Kennungen REQ-82 und REQ-83 und übernahm die Claim-Referenz samt Herkunftsangaben des Eingangs --- das zugehörige Feature und seine Arbeitspakete entstehen erst danach. \\
    Arbeitspakete und Themenzuordnung & Eine nachgelagerte modellgestützte Stufe schlug ein neues Feature („Besuchsankündigung und Besuchskoordination``) mit fachlicher Begründung vor; nach eigener menschlicher Freigabe (in diesem Lauf als dokumentierte Sammelfreigabe im UI) legte die Anwendung FC-15 und PBI-046/047 an, verband die PBIs über „deckt`` mit ihren Anforderungen und ordnete PBIs sowie Anforderungen dem Feature zu. \\
    Ergebnis & Die beiden Arbeitspakete wurden als GitHub-Issues \#44 („Besuche \ldots{} vorab ankündigen``) und \#45 („Übersicht der angekündigten Besuche \ldots``) real erzeugt; die Zuordnung PBI-zu-Issue wurde im Core gespeichert. \\
    \hline
  \end{tabular}
  \caption[Besuchsankündigung: von der Quelle zum Issue]{Ein
    fachlicher Inhalt über die Stationen des ausgeführten Falls F1
    (Abschnitt~\ref{sec:eval-fortschreibung}; Belegzuordnung in
    Anhang~\ref{anh:evaluation}). Die Quellenzeile folgt dem
    Originaltranskript, die übrigen Zeilen den Lauf-Artefakten.}
  \label{t:beispiel-besuchsplanung}
\end{table}


Die Stationen verändern jeweils etwas anderes: Die modellgestützte
Ableitung zerlegt den gebündelten Claim in zwei normativ formulierte
Anforderungen und erhält den Claim-Bezug als Herkunftsangabe. Die
menschliche Auswahl entscheidet anschließend über die Übernahme
(Abschnitt~\ref{sec:governance-mutation}). Erst die nachgelagerte,
eigens freigegebene Arbeitspaket-Stufe ergänzt Feature-Zuordnung und
„deckt``-Beziehungen. Ihre Freigabe erfolgte in diesem Lauf als
Sammelfreigabe; daraus lässt sich keine Aussage über die Sorgfalt
einer Einzelprüfung ableiten. Die Projektionsbeziehung
(„umgesetzt durch Issue``) verbindet die Arbeitspakete schließlich
mit Issues~\#44/\#45. Sie bezeichnet hier ihre operative
Darstellung als Arbeitsaufträge, keine bereits implementierte
Softwarefunktion.

Die Quellenzeile der Tabelle folgt dem Originaltranskript. Im
gespeicherten Ledger-Beleg ist der Besuchsaussage dagegen das
abweichende Sprecherpräfix „Leitung:`` vorangestellt. Die nachträgliche
Code- und Artefaktprüfung zeigte, dass der Parser die Rollenbezeichnung
„Angehörigen-Vertreterin:`` nicht als Sprecherwechsel erkannt hatte.
Die am 11.09.2026 gesondert geprüfte Korrektur verändert den
historischen Beleg nicht (Anhang~\ref{anh:evaluation}). Seine
Zuschreibung wird nicht als wörtlich identisch mit der Quelle
behandelt. Der Fall veranschaulicht damit auch, weshalb eine
vorhandene Herkunftsreferenz die Prüfung ihres Inhalts nicht ersetzt.

\emph{Derselbe Sachverhalt aus einem anderen Eingang.}
Ein späterer Kommentar zu Issue~\#45 fordert, die Besuche der nächsten
14 Tage anzuzeigen. Zu diesem Zeitpunkt bestehen die allgemeine
Übersichtsanforderung REQ-83, das Feature FC-15 und das Arbeitspaket
PBI-047 mit seiner Issue-Zuordnung bereits im Core.
Abbildung~\ref{fig:core-netz-besuchsbeispiel} vertieft daran den
Betriebspfad der Gesamtarchitektur aus
Abbildung~\ref{fig:logische-gesamtarchitektur}: Sie zeigt, wie die
zusätzliche Information übernommen und mit dem bestehenden
Projektbestand verbunden wird. Links stehen sechs fachliche
Stationen in Verarbeitungsreihenfolge; rechts sind Ausgangsbestand,
Ergebnis nach Schritt 5 und Issue-Aktualisierung nach Schritt 6
dargestellt.

Die Rückmeldung gelangt über die \emph{GitHub-Ernte} in das System
(Abschnitt~\ref{sec:projektionen-kontrollierte-aussenkopplung}).
Dieser Rückweg wird über den Direktzugang oder den Steward gestartet.
Auf Abruf und Sammlung verarbeitbarer Änderungen aus Issues und
Kommentaren folgt die fachliche Aufbereitung durch Agenten.
Im Beispiel wird aus dem Kommentar ein Vorschlag; die deterministische
Paketbildung übernimmt ihn mit seinen Herkunftsangaben in ein
\emph{Delta}. Schritt 1 fasst diese Aufbereitung zusammen.
Das Delta enthält zu prüfende Inhalte, noch keine autorisierten
Änderungen.

Ab der Einordnung gegen den vorhandenen Core (Schritt 2) nutzt dieser
Eingang denselben Fortschreibungsablauf wie die aufbereiteten Inhalte
aus Folgemeetings oder dialogischen Eingaben
(Abschnitt~\ref{sec:logische-gesamtarchitektur}). Beim Meeting führen
Ledger-Verarbeitung und fachliche Ableitung zum Delta, bei der
dialogischen Eingabe dessen kontrollierte Aufbereitung über den
Steward. Dieser dient zugleich als Bedienebene für verschiedene
Abläufe. Welche Prüfungen, Freigaben und Folgeoperationen benötigt
werden, hängt von Inhalt und Einordnung ab.

Im dargestellten Fall schlägt der Einordnungsagent nach Abgleich mit
dem Bestand eine zusätzliche Anforderung vor. Die über den Steward
eingereichte menschliche Freigabe erlaubt deren Übernahme. In
Schritt 3 legt die deterministische Anwendung REQ-89 mit den
Herkunftsangaben des Eingangs an. Die ergänzende Anforderung betrifft
den 14-Tage-Zeitraum; REQ-83 und das Besuchsthema bleiben bestehen.
Die Einordnung könnte bei anderen Inhalten auch zu einer Bestätigung,
Änderung oder einem Klärungsbedarf führen
(Abschnitt~\ref{sec:core-projektfortschreibung}).

Die Arbeitspaket-Stufe kann bestehende PBIs und Feature-Zuordnungen
ergänzen oder neue Strukturen vorschlagen. Die Schritte 4 und 5
betreffen hier die bereits geplante Arbeit. Agenten
schlagen vor, PBI-047 um REQ-89 zu erweitern und seinen Inhalt
anzugleichen. Nach eigener Freigabe entstehen die zusätzliche
\emph{deckt}-Beziehung und die Zuordnung von REQ-89 zum bestehenden
Feature FC-15. Titel und Akzeptanzkriterien berücksichtigen nun den
14-Tage-Zeitraum. Schritt 6 aktualisiert nach Freigabe der
Veröffentlichung das bereits zugeordnete Issue~\#45.

\begin{figure}[htbp]
  \centering
  \includegraphics[width=\textwidth,height=0.84\textheight,keepaspectratio]{figures/05/core-netz-besuchsbeispiel.pdf}
  \caption[Schrittweise Fortschreibung am Besuchsbeispiel]{
    Fortschreibung der Besuchsübersicht durch einen GitHub-Kommentar.
    Der Themenrahmen zeigt die Feature-Zuordnung; die Pfeile vom
    Arbeitspaket zu den Anforderungen bezeichnen deren geplante
    Umsetzung. Der gestrichelte Pfeil zeigt den Lesezugriff auf den
    Core, H die menschlichen Freigaben. Quellenangaben fassen
    archivierte Belegwege zusammen (Anhang~\ref{anh:provenienz-v4};
    eigene Darstellung, Bezeichnungen gekürzt).}
  \label{fig:core-netz-besuchsbeispiel}
\end{figure}

% P10-6-Nachzug, 16.09.2026: Erklärabsatz nicht durch eine ganzseitige Übersicht trennen.
\begin{samepage}
Die als „Quelle`` bezeichneten Angaben benennen den ursprünglichen
fachlichen Beleg. Für REQ-83 führt der Herkunftsweg über den Ledger
zum Meeting-Transkript. Für REQ-89 führen die Angaben zu Eingang
und Erzeugungslauf zum archivierten Eingangspaket; dort ist der
konkrete Kommentar zu Issue~\#45 zugeordnet. Der Quellenhinweis
fasst diesen Belegweg zusammen. Kommentar und spätere Aktualisierung
des Issue-Inhalts haben unterschiedliche Funktionen: Der Kommentar
liefert die zusätzliche Vorgabe, die Projektion stellt die
freigegebene Arbeit operativ dar.
\par
\end{samepage}

% P10-6: Erklärung der Erzeugerangabe nicht mit einer einzelnen Zeile vor der Grafik beginnen.
\begin{samepage}
Die in Abschnitt~\ref{sec:core-projektfortschreibung} eingeführte
Erzeugerangabe unterscheidet hier Beziehungen aus der ursprünglichen
Arbeitspaketbildung von denen der späteren PBI-Fortschreibung.
Die zeitliche Zuordnung stützt sich zusätzlich auf archivierte
Laufstände, Pläne, Entscheidungen und Übernahmeberichte. Die Strukturprüfung vor
dem Speichern kontrolliert die definierten Integritätsregeln
(Abschnitt~\ref{sec:governance-mutation}). Die fachliche Richtigkeit
der Zuordnung bleibt davon zu unterscheiden. Der Fall macht damit
die Aufgabenteilung sichtbar: Agenten interpretieren und schlagen
vor, der Workflow führt durch die vorgesehenen Prüfungen und
Freigaben, und die deterministische Anwendung setzt die
freigegebenen Inhalte und Verbindungen um.
\par
\end{samepage}

Andere Zustandswirkungen benötigen eigene Belege: F2 betrifft die
Verfeinerung der Medikamenteneinsicht bei erhaltener Kennung und
Historie (Realisierung in
Abschnitt~\ref{sec:realisierung-agent-executor-workflow},
Fallprüfung in Abschnitt~\ref{sec:eval-fortschreibung}).
E47-02 illustriert den Konflikt- und Angleichungspfad unter den in
Abschnitt~\ref{sec:einordnende-fortschreibung} genannten
Experimentbedingungen. E48-03 dokumentiert separat den Weg einer
Außenrückmeldung zum kontrollierten Eingang. Diese Fälle ergänzen
das Besuchsbeispiel, ohne zu demselben Lauf zu gehören.

Die beiden Pfade nutzen dieselbe Architektur, unterscheiden sich
aber beim ersten Schreibvorgang: Für den initialen Seed fehlt die
Einzelautorisierung, während die beschriebenen Fortschreibungen
des vorhandenen Bestands den regulären Autorisierungspfad nutzen.
Abschnitt~\ref{sec:architektursynthese-gestaltungsfolgen} fasst die
daraus folgenden Gestaltungsfolgen und Grenzen zusammen.
